
\begingroup
\setlength{\parskip}{0.2em}

\chapter*{符号清单与术语表}\addcontentsline{toc}{chapter}{符号清单与术语表}

本文中频繁使用的主要符号、核心术语与常用缩写汇总如下。为避免与正文重复，本页重点列出第三章方法定义与后续评价分析中反复出现的记号与表述。

{\noindent\large\textbf{符号清单}\par}
\vspace{0.2em}

\noindent\textbf{基本对象与集合}

\begingroup
\small
\centering
\setlength{\tabcolsep}{3pt}
\renewcommand{\arraystretch}{1.03}
\begin{tabularx}{0.90\textwidth}{>{\centering\arraybackslash}p{2.9cm} >{\raggedright\arraybackslash}X}
\toprule
\textbf{符号} & \textbf{含义说明} \\
\midrule
$\mathcal{V}$ & 漏洞实例空间，表示全部待分析漏洞构成的集合。 \\
$v=\langle \mathit{id},d,m \rangle$ & 单个漏洞实例，由编号、描述与元数据组成。 \\
$\mathit{id}$ & 漏洞标识符，如 CVE 编号。 \\
$d$ & 漏洞描述文本，用于刻画漏洞事实。 \\
$m$ & 漏洞元数据，如严重性、受影响组件与版本范围等。 \\
$\mathcal{C}$ & 安全上下文空间，表示与漏洞相关的证据集合空间。 \\
$c$ & 安全上下文，即带来源标识的有限证据集合。 \\
$t_i$ & 第 $i$ 条证据内容，如代码片段、公告文本或文档片段。 \\
$\mathit{src}_i$ & 第 $i$ 条证据来源，如文档路径、链接或代码位置。 \\
$\Pi$ & 候选策略全集，表示所有候选监控策略与防护策略的总集合。 \\
$\pi=\langle \mathit{type},\mathit{method},\mathit{action} \rangle$ & 单个安全策略，由策略类型、控制机理与执行表达组成。 \\
$\mathit{type}$ & 策略类型，取值为监控型或防护型。 \\
$\mathit{method}$ & 控制机理序列，用于表示策略采用的控制手段。 \\
$\mathit{action}$ & 执行表达序列，对应可部署的规则、配置或命令。 \\
$\Pi_v$ & 漏洞 $v$ 的候选策略集合，由生成阶段针对漏洞 $v$ 产出。 \\
$\mathcal{T}=\{\mathrm{mon},\mathrm{prot}\}$ & 策略类型集合，分别表示监控型与防护型。 \\
$\tau$ & 类型变量，表示某一给定策略类型。 \\
$\Pi_v^{\tau}$ & 同类型候选子集，表示漏洞 $v$ 在类型 $\tau$ 下的候选集合。 \\
\bottomrule
\end{tabularx}
\endgroup

\vspace{0.35em}

\noindent\textbf{评价与聚合记号}

\begingroup
\small
\centering
\setlength{\tabcolsep}{3pt}
\renewcommand{\arraystretch}{1.03}
\begin{tabularx}{0.90\textwidth}{>{\centering\arraybackslash}p{2.9cm} >{\raggedright\arraybackslash}X}
\toprule
\textbf{符号} & \textbf{含义说明} \\
\midrule
$\Delta=\{E,I,D\}$ & 评价维度集合，分别对应有效性、无干扰性与可部署性。 \\
$q(\pi)$ & 三维质量向量，表示策略 $\pi$ 在三个维度上的质量。 \\
$E(\pi)$ & 有效性，用于衡量对主要攻击路径与关键风险面的覆盖程度。 \\
$I(\pi)$ & 无干扰性，衡量对正常业务的扰动程度及误报风险。 \\
$D(\pi)$ & 可部署性，衡量策略在现有环境中的落地难度。 \\
$\mathcal{H}$ & 评价主体集合，表示多个评审主体构成的集合。 \\
$r_{\delta}^{(h)}$ & 秩函数，表示评价主体 $h$ 在维度 $\delta$ 上给出的排序位置。 \\
$s_{\delta}^{(h)}(\pi)$ & 归一化秩得分，用于将排序位置归一化到 $[0,1]$ 区间。 \\
$\hat S_{\delta}(\pi)$ & 共识得分，表示多个评价主体在维度 $\delta$ 上的聚合结果。 \\
$U_{\delta}(\pi)$ & 分歧度量，用于刻画不同评审结果的一致性程度。 \\
$Q(\pi)$ & 总体质量函数，表示综合三维共识得分后的最终质量。 \\
$\alpha_E,\alpha_I,\alpha_D$ & 维度权重，分别对应有效性、无干扰性与可部署性的权重。 \\
\bottomrule
\end{tabularx}
\endgroup

\vspace{0.45em}

{\noindent\large\textbf{术语表}\par}
\vspace{0.2em}

为统一全文表述，本文使用频率较高的核心术语与常用缩写说明如下。

\noindent\textbf{核心术语}

\begingroup
\small
\centering
\setlength{\tabcolsep}{3pt}
\renewcommand{\arraystretch}{1.03}
\begin{tabularx}{0.90\textwidth}{>{\centering\arraybackslash}p{1.8cm} >{\centering\arraybackslash}p{3.0cm} >{\raggedright\arraybackslash}X}
\toprule
\textbf{术语} & \textbf{英文/缩写} & \textbf{说明} \\
\midrule
云原生 & Cloud-native & 本文研究对象所处的应用与基础设施环境。 \\
Kubernetes & K8s & 本文默认指 Kubernetes 容器编排平台，K8s 为其常用缩写。 \\
权限提升 & Privilege Escalation & 攻击者突破原有权限边界并获取更高权限的过程。 \\
补丁窗口期 & Patch-gap window & 漏洞披露至补丁可用或补丁完成部署之间的时间段。 \\
临时防控 & Interim mitigation & 在补丁窗口期内用于降低风险的临时控制措施。 \\
安全策略 & Security strategy & 监控策略与防护策略的上位概念。 \\
监控策略 & Monitoring strategy & 面向异常行为感知、告警与取证的策略。 \\
防护策略 & Protection strategy & 面向攻击阻断、权限收敛与攻击面控制的策略。 \\
安全上下文 & Security context & 与具体漏洞相关的证据集合及其来源标识。 \\
安全知识库 & Security knowledge base & 存放多源证据、向量索引及元数据的知识支撑体系。 \\
多智能体协同 & Multi-agent collaboration & 由多个功能角色协同完成分析、生成与评审的工作方式。 \\
检索增强 & RAG & Retrieval-Augmented Generation \\
思维链推理 & CoT & Chain-of-Thought，思维链。 \\
反馈优化 & Feedback & 基于评审结果对候选策略进行修订与收敛的过程。 \\
\bottomrule
\end{tabularx}
\endgroup

\vspace{0.35em}

\noindent\textbf{常用缩写}

\begingroup
\small
\centering
\setlength{\tabcolsep}{3pt}
\renewcommand{\arraystretch}{1.03}
\begin{tabularx}{0.90\textwidth}{>{\centering\arraybackslash}p{2cm} >{\raggedright\arraybackslash}p{4.9cm} >{\raggedright\arraybackslash}X}
\toprule
\textbf{缩写} & \textbf{英文全称} & \textbf{中文说明} \\
\midrule
CVE & Common Vulnerabilities and Exposures & 通用漏洞与暴露。 \\
CWE & Common Weakness Enumeration & 通用弱点枚举。 \\
CVSS & Common Vulnerability Scoring System & 通用漏洞评分系统。 \\
NVD & National Vulnerability Database & 国家漏洞数据库。 \\
CNCF & Cloud Native Computing Foundation & 云原生计算基金会。 \\
K8s & Kubernetes & Kubernetes 的常用缩写。 \\
RBAC & Role-Based Access Control & 基于角色的访问控制。 \\
LLM & Large Language Model & 大语言模型。 \\
\bottomrule
\end{tabularx}
\endgroup

\endgroup
